%% ======================================================================
%% chap4/methodology_main.tex
%%
%% Chapter 4: Methodology and System Design
%%
%% Presents the research methodology and the resulting system design.
%% Section 4.1 covers design science research framing, iterative
%% prototyping, data collection, and evaluation philosophy; Section 4.2
%% covers the conceptual system design (unified rights token, three
%% monetization models, NFT nesting, cross-chain model, pallet-first
%% architecture, layered architecture).
%%
%% Sections:
%%   4.1 Research Methodology (5 subsections)
%%   4.2 System Design (6 subsections)
%%
%% External labels referenced:
%%   chap:background (Ch 2), chap:implementation (Ch 5),
%%   chap:evaluation (Ch 6), chap:results (Ch 7),
%%   sec:pallet-impl (5.2), sec:cross-chain-integration (5.3),
%%   sec:challenges-solutions (5.4), sec:kpis (6.5),
%%   sec:challenges-and-gaps (2.4),
%%   app:work-plan (Appendix F),
%%   sec:rmrk-to-pallet-nfts (Appendix A.1),
%%   sec:contract-first-pivot (Appendix A.2),
%%   sec:children-bound (Appendix A.8),
%%   sec:ink-discontinuation (Appendix A.9),
%%   subsec:bounded-vec-future (9.3.4 in Table 9.1)
%%
%% Citations: hevnerEtAl2004 (the canonical DSR reference)
%% ======================================================================

\chapter{Methodology and System Design}\label{chap:methodology}

This chapter describes the research methodology we used to develop CCRMS and the resulting system design. Section~\ref{sec:research-methodology} covers the methodology, encompassing design science framing, iterative prototyping, data collection, and evaluation philosophy. Section~\ref{sec:system-design} describes the system design at the conceptual level; Chapter~\ref{chap:implementation} covers implementation details.

\section{Research Methodology}\label{sec:research-methodology}

\subsection{Design Science Research Framing}\label{subsec:dsr-framing}

CCRMS is fundamentally a systems-building project, with a primary artifact: a working software prototype. The most appropriate methodological framework is therefore \textbf{design science research (DSR)}, as defined by \citet{hevnerEtAl2004}. Our dissertation maps to Hevner's seven DSR guidelines as follows: (1) \emph{design as an artifact:} the CCRMS parachain and seven implementation specifications; (2) \emph{problem relevance:} the commercial and academic gaps identified in Chapter~\ref{chap:background}; (3) \emph{design evaluation:} eleven KPIs, covered in Chapters~\ref{chap:evaluation}--\ref{chap:results}; (4) \emph{research contributions:} the artifact, the design knowledge captured in the specifications, and the empirical measurements; (5) \emph{research rigor:} established Polkadot SDK primitives, standard Rust testing, and peer-reviewed Monte Carlo methodology; (6) \emph{design as a search process:} nine significant pivots documented in Section~\ref{sec:challenges-solutions}; (7) \emph{communication of research:} this dissertation and an open-source code release.

DSR is especially appropriate here because it recognizes both the artifact and the design knowledge it generates as coequal contributions.

\subsection{Iterative Prototyping}\label{subsec:iterative-prototyping}

Within the DSR framework, the development process follows an \textbf{iterative prototyping} model (build, test, measure, refine). We chose iteration deliberately: the Polkadot SDK monorepo releases stable versions approximately quarterly, rendering a waterfall approach susceptible to conclusions based on an outdated toolchain, and the architectural intricacies of a cross-chain system imply that certain design decisions can only be thoroughly assessed once the system is operational. In practice, this iterative methodology resulted in nine substantial architectural pivots during Phase~4, which are documented as challenges in Section~\ref{sec:challenges-solutions}. We executed the work across six four-week phases (Appendix~\ref{app:work-plan}); references to ``Phase~N'' throughout the dissertation denote work produced in that phase.

\subsection{Data Collection Methods}\label{subsec:data-collection}

Our evaluation employs a mixed-methods approach: quantitative measurements such as unit tests, throughput and latency benchmarks, storage growth, cross-chain XCM latency, reliability tests, a comparative centralized benchmark, and a Monte Carlo creator revenue simulation. Additionally, qualitative inputs are derived from a literature synthesis. Chapter~\ref{chap:evaluation} describes each method. All quantitative outputs are systematically stored as machine-readable JSON files in the \texttt{scripts/perf/results/} directory, ensuring reproducibility.

\subsection[Evaluation Philosophy: Measurement-Focused]{Evaluation Philosophy: Measurement-Focused, Not Target-Driven}\label{subsec:evaluation-philosophy}

We adopt a \textbf{measurement-focused philosophy}: the goal is to document and analyze the levels of performance, efficiency, security, and economic viability achieved by the system, rather than to verify that it conforms to predetermined targets. This measurement-centric approach maintains academic rigor regardless of whether specific results meet hypothetical benchmarks, renders unexpected results valuable rather than embarrassing, encourages truthful reporting, and establishes a baseline for future researchers.

Each of the eleven KPIs in Section~\ref{sec:kpis} carries a target drawn from the Phase~2 methodology, but we treat those targets as reference points rather than pass/fail criteria. In practice, all eleven met their targets; the methodology would have remained valid if some had not.

\subsection{Methodology Pivots}\label{subsec:methodology-pivots}

We revised the methodology established in Phase~4. The original plan called for Kusama via the Canary Chaos program, implementing RMRK 2.0 composable NFTs, scheduling XCM for automatic renewal, and conducting a Monte Carlo simulation, which had not yet been executed. The actual implementation modified this approach by substituting Kusama with local Zombienet to facilitate faster iteration, replacing RMRK 2.0 with \texttt{pallet-nfts} (noting that the RMRK pallets are frozen on SDK version 0.9.36), replacing the scheduled XCM with an \texttt{on\_initialize} block hook (considering that the XCM v5 \texttt{Schedule} instruction is not yet production-ready), and completing the Monte Carlo simulation in Phase~5.

The most significant methodological pivot is unrelated to the original plan: the \textbf{emphasis on architectural decisions and pivots as first-class research outcomes}. The Phase~2 methodology regarded implementation merely as a means to generate measurements; however, in practice, the architectural decisions made during the pivots emerged as equally valuable contributions as the measurements themselves.

\section{System Design}\label{sec:system-design}

This section outlines the conceptual design. Section~\ref{sec:pallet-impl} provides the implementation details.

\subsection{The Unified Rights Token Concept}\label{subsec:unified-rights-token}

The principal intellectual contribution of CCRMS is the \textbf{unified rights token}: a single on-chain asset that can represent three distinct modes of content access (an active subscription, a remaining pay-per-view balance, or permanent ownership) for a single piece of content from a single creator under a unified set of rules.

The unified token addresses the gap identified as G7 in Section~\ref{sec:challenges-and-gaps}. Current blockchain platforms compel creators and consumers to choose among various monetization models: subscription platforms cannot represent per-view or ownership rights, NFT marketplaces cannot facilitate recurring access, and pay-per-view systems seldom function effectively on-chain. The unified token treats the three models as \textbf{access modes} of a single underlying content registration, rather than categorizing them as separate platforms. A creator registers content once and can offer it through all three modes concurrently; a consumer chooses the mode that best suits their needs; the system enforces all three from a single canonical record. This approach eliminates a source of friction in the market: creators no longer need to maintain three distinct platform accounts for different audiences, and consumers no longer have to navigate three separate interfaces with separate payment flows.

\subsection{The Three Monetization Models}\label{subsec:three-models}

The architecture supports three models, each motivated by a different gap in the existing landscape:

\begin{itemize}
\item \textbf{Subscription:} time-limited access valid for a configurable number of blocks per content item, with optional automatic renewal driven by an on-chain block hook.
\item \textbf{Pay-per-view (PPV):} prepaid access to a specified number of views of a content item is available. Each \texttt{consume\_view} operation reduces the counter; additional purchases increase the count rather than overwrite it. Conceptualizing pay-per-view (PPV) as a counter, rather than a sequence of individual view transactions, circumvents the micro-transaction issue that has historically impeded on-chain PPV models.
\item \textbf{Permanent ownership:} unrestricted and transferable access is granted. Ownership is prioritized during the access verification process. Transfers adhere to a mint-burn-mint pattern rather than a direct NFT transfer, thereby safeguarding provenance and attributive consistency.
\end{itemize}

A single content item can offer all three models concurrently with independently configured pricing. The \texttt{check\_access} function resolves rights in priority order (Ownership~>~Subscription~>~PPV).

\subsection{NFT-Based Rights Tokenization}\label{subsec:nft-tokenization}

CCRMS encodes content rights as non-fungible tokens using a parent--child nesting model atop \texttt{pallet-nfts}. The decision to use NFTs rather than a bespoke storage configuration was motivated by their widespread acceptance as a standard, their compatibility with existing digital wallets, their capability for independent transferability, which facilitates a secondary market for ownership rights, and the clear separation of the content item's identity (the parent) from user-specific rights derived from it (the children).

The conceptual model: each content registration is an NFT collection comprising one parent Content NFT and zero or more child rights NFTs, each tagged with a \texttt{rights\_type} attribute (Subscription, PPV, or Ownership) and owned by the respective user. The parent--child structure fulfills three concurrent purposes: establishing identity, tokenizing individual rights, and enabling on-chain auditability of the rights-to-content relationship.

The nesting pattern is inspired by RMRK 2.0, the composable NFT specification for Polkadot; however, it is implemented directly on top of \texttt{pallet-nfts} rather than depending on the \texttt{rmrk-substrate} pallets, which are fixed on SDK 0.9.36. Appendix~\ref{sec:rmrk-to-pallet-nfts} provides a detailed account of this transition. The outcome illustrates that RMRK-style composability can be achieved solely through standard FRAME primitives.

One significant constraint is that the \texttt{Children} storage entry is a \texttt{BoundedVec<\_, ConstU32<50>{}>}, capping child NFTs to fifty per parent. This restriction is necessary for FRAME benchmarking, as the worst-case weight must be calculable in advance. The examination of this constraint is detailed in Appendix~\ref{sec:children-bound} and treats mitigation as future work in Section~\ref{subsec:bounded-vec-future}.

\subsection{Cross-Chain Communication Model}\label{subsec:cross-chain-model}

Cross-chain operation is the mechanism by which the unified rights token model addresses portable digital ownership: a user who acquires rights on one blockchain should be able to use them on another without forfeiting access or repurchasing. The model possesses two conceptual properties.

\textbf{Payer--beneficiary decoupling.} Each cross-chain rights operation distinguishes between the funding account and the rights recipient account. In a local call, the payer and the beneficiary are identical (the caller); in a cross-chain call, the dispatch origin is the sovereign account of the remote parachain, and the beneficiary is a distinct user account. This configuration enables a remote parachain to acquire rights on behalf of its users without necessitating their presence on CCRMS.

\textbf{Trustless verification of remote rights.} A second pallet, \texttt{pallet-rights-verifier}, provides Merkle proof verification of the CCRMS rights state to other parachains. A consuming chain can verify, without relying on any oracle or relayer, that a specific account presently maintains an active subscription by submitting a storage proof against a known CCRMS state root. This approach upholds the principle that CCRMS rights are \textbf{canonical on the CCRMS chain}: other chains authenticate cryptographically rather than maintaining replicas.

The model also extends beyond Polkadot via \textbf{Snowbridge}, a trust-minimized bridge connecting Polkadot and Ethereum. We configure the runtime to accept messages originating from Ethereum via the standard \texttt{LocationToAccountId} converters. Our architectural choice at this stage was to guarantee \textbf{forward compatibility with Ethereum bridging without reliance on it}. Section~\ref{sec:cross-chain-integration} provides a comprehensive account of the complete cross-chain integration.

\subsection{Pallet-First Architecture}\label{subsec:pallet-first}

The most significant architectural decision was to centralize the business logic within a \textbf{FRAME pallet} rather than in smart contracts. This decision was made during Phase~4 for three reasons: the operational maturity disparity between FRAME pallets and \texttt{pallet-revive}; the unavailability of the chain extension mechanism, which the original contract-first approach relied upon, for write operations; and a pallet-first design with standard Polkadot wallets without EVM-style adaptation. Appendix~\ref{sec:contract-first-pivot} provides a detailed documentation of this decision, while Appendix~\ref{sec:ink-discontinuation} notes how it was unexpectedly validated by the discontinuation of ink! development in January 2026. The overarching principle, that architectural choices which reduce dependency on any non-essential component tend to yield benefits during ecosystem-level changes, is reflected in several other decisions within CCRMS.

\subsection{Layered Architecture}\label{subsec:layered-architecture}

The system is organized as six distinct, independently replaceable layers: (1) the \textbf{Polkadot relay chain} offers shared security and facilitates HRMP message passing; (2) the \textbf{parachain runtime} composes the CCRMS FRAME pallets, configures the XCM executor, and incorporates a custom EVM precompile; (3) the \textbf{business logic} layer comprises \texttt{pallet-content-rights} (17 extrinsics, 10 storage items) and \texttt{pallet-rights-verifier}; (4) the \textbf{NFT layer} is the standard \texttt{pallet-nfts}, with parent--child nesting enabled by \texttt{pallet-content-rights}; (5) the \textbf{smart contract~/ precompile layer} provides EVM- and ink!-compatible access through the \texttt{rights\_manager} contract and the \texttt{ContentRightsPrecompile}, both of which are non-load-bearing wrappers; (6) the \textbf{cross-chain layer} manages XCM v3 messaging in addition to the Snowbridge integration. Independence among layers is a fundamental design objective: the smart contract layer may be removed or replaced, the XCM protocol may evolve from v3 to v5, and the relay chain may transition to JAM, all without affecting the pallet code.
